\subsubsection{Bình phương tối thiểu thông thường (Ordinary least squares,
OLS)}

Phương pháp Bình phương tối thiểu thông thường (OLS) hướng đến việc tìm bộ
trọng số $\boldsymbol{w}^\ast$ sao cho cực tiểu quá tổng bình phương sai số giữa giá trị
thực tế và giá trị dự đoán hay cực tiểu hoá hàm lỗi: \begin{equation}
  E_D(\boldsymbol{w}) = 1/2\cdot\lVert \boldsymbol{t} -
  \boldsymbol{\Phi}\boldsymbol{w}\rVert^2.
\end{equation}

\paragraph{Phương trình chuẩn (Normal Equation)}

Phương trình chuẩn cung cấp một lời giải đóng cho bài toán tối ưu trên bằng
cách giải nghiệm đạo hàm bậc nhất bằng 0. Công thức nghiệm tối ưu là:
\begin{equation}
  \boldsymbol{w}^\ast = \arg\min_{\boldsymbol{w}}E_D(\boldsymbol{w}) =
  (\boldsymbol{\Phi}^\intercal\boldsymbol{\Phi})^{-1}
  \boldsymbol{\Phi}^\intercal\boldsymbol{t}.
\end{equation}

Trong thực tế cài đặt, thay vì nghịch đảo trực tiếp
($\boldsymbol{\Phi}^\intercal\boldsymbol{\Phi})^{-1}$---vốn dễ gây lỗi nếu
ma trận bị suy biến (singular)---em sử dụng ma trận giả nghịch đảo
Moore--Penrose $\boldsymbol{\Phi}^\dagger$. Công thức triển khai trở thành:
\begin{equation}
\boldsymbol{w}^\ast = \boldsymbol{\Phi}^\dagger \boldsymbol{t}.
\end{equation}

Độ phức tạp tính toán nghiệm là $\Theta(NM^2 + NM + M^3 + M^2)$. Đối với bộ dữ
liệu California Housing, khi mà $N \gg M$ ($20640 \gg 9$, trong đó $9$ là số
lượng đặc trưng cộng thêm 1 để học hệ số chặn), độ phức tạp tương
đương $\Theta(NM^2)$.

\paragraph{Mini-batch Gradient Descent}

Đối với các bài toán có kích thước dữ liệu lớn, việc tính toán trực tiếp trên
toàn bộ dữ liệu gây áp lực lớn cho bộ nhớ. Thuật toán Mini-batch
Gradient Descent nhằm cân bằng giữa tốc độ hội tụ và hiệu năng tính toán.

Về hiệu suất, độ phức tạp tính toán của một lần xử lý một lô (batch) là
$\Theta(BM)$. Khi đó, với một epoch, độ phức tạp là $\Theta(NM)$. Và với $k$
epoch thì độ phức tạp là $\Theta(kNM)$.

\paragraph{So sánh giữa phương pháp Normal Equation và Gradient Descent}
